% ===================================================================
%  Front-tracking and VOF for THREE-DIMENSIONAL two-phase flow
%  Newtonian and viscoelastic, on uniform and adaptive meshes
% ===================================================================
%
% PARA O OVERLEAF: suba esta pasta inteira (main.tex, preambulo.tex,
% secoes/, dados/, referencias.bib) e marque main.tex como documento
% principal.  Compilador: pdfLaTeX.  Nao ha' nenhuma imagem binaria -- todos
% os graficos saem do pgfplots lendo dados/*.dat --, entao a pasta e' pequena
% e autocontida.
%
% Localmente: latexmk -pdf main   (exige texlive-publishers, pelo elsarticle)
%
% O DESENHO E' UM 2x2x2 e a estrutura segue dele: dois METODOS
% (front-tracking, VOF) x duas REOLOGIAS (newtoniana, viscoelastica) x duas
% MALHAS (uniforme, adaptativa), tudo em TRES DIMENSOES.  As oito celulas
% dessa matriz sao as oito corridas de producao que o artigo tem de sustentar.
%
% ESTADO DE CADA SECAO -- o que ja' tem respaldo medido e o que nao tem.
% Cada secao abre com um bloco \status{} dizendo isso, e esses blocos somem do
% PDF quando \nostatus e' definido.  Existem para que ninguem -- nos inclusive
% -- leia um paragrafo de planejamento como resultado.
%
%   1 Introduction           a escrever
%   2 Formulation            PRONTA: equacoes e convencao medidas
%   3 Numerical method       PRONTA: todas as partes implementadas
%   4 Verification           PRONTA: doze oraculos medidos
%   5 Newtonian 3D           FT vs VOF uniforme FECHADO em N=25 e N=50;
%                            caso 2 nao rodado; adaptativo nao em producao
%   6 Viscoelastic 3D        matriz uniforme FECHADA (12 corridas);
%                            cuspide na resolucao de referencia FECHADA
%   7 Cost and performance   taxas em malha uniforme MEDIDAS
%   8 Conclusions            depende das acima

\documentclass[preprint,review,12pt]{elsarticle}
\journal{Journal of Computational Physics}

% ---------------------------------------------------------------- preambulo
% Comum aos dois artigos (2D e 3D).  Mantido num arquivo so' para que uma
% mudanca de estilo nao precise ser feita duas vezes e divergir na segunda.
\usepackage[T1]{fontenc}
\usepackage[utf8]{inputenc}
\usepackage[english]{babel}
\usepackage{lmodern}
\usepackage{microtype}
\usepackage{amsmath,amssymb}
\usepackage{booktabs}
\usepackage{adjustbox}
\usepackage{siunitx}
\usepackage{graphicx}
\usepackage{pgfplots}
\usepackage{pgfplotstable}
\usepackage{tikz}
\usepackage{xcolor}
\pgfplotsset{compat=1.18}

% OS NUMEROS VIVEM EM dados/*.dat e sao lidos por pgfplotstable, nunca
% digitados no texto.  Regerar uma medida troca o .dat e o texto acompanha,
% sem chance de a tabela e a prosa discordarem.
\newcommand{\dat}[1]{dados/#1}

% \status{...}: o que a secao tem de respaldo medido.  Definir \nostatus antes
% do \begin{document} faz os blocos sumirem na versao de submissao.
\newcommand{\status}[1]{%
  \ifdefined\nostatus\else
    \begingroup
    \par\medskip\noindent
    \colorbox{yellow!18}{\begin{minipage}{0.97\linewidth}
      \small\textbf{Status of this section.} #1
    \end{minipage}}\par\medskip
    \endgroup
  \fi}

% \missing{...}: uma lacuna pontual, visivel de proposito.
\newcommand{\missing}[1]{\textcolor{red!70!black}{\textbf{[missing:} #1\textbf{]}}}

% \caveat NAO e' o mesmo que \missing.  \missing marca o que falta MEDIR;
% \caveat marca o que foi medido mas com uma ressalva que muda a leitura.
% Confundir os dois esconderia um resultado-com-reserva entre pendencias.
\newcommand{\caveat}[1]{\par\smallskip\noindent\textcolor{orange!60!black}{%
  \textbf{[caveat:} #1\textbf{]}}\par\smallskip}

\usepackage[hidelinks]{hyperref}

% Descomente para a versao de submissao, sem os blocos de estado:
% \def\nostatus{}

\emergencystretch=3em

\begin{document}

\begin{frontmatter}

\title{Front-tracking and VOF for three-dimensional two-phase flow:
       Newtonian and viscoelastic, on uniform and adaptive meshes}

\author[icmc]{Antonio Castelo Filho}
\ead{castelo@icmc.usp.br}
\author[icmc]{\missing{co-authors}}

\affiliation[icmc]{organization={Instituto de Ci\^encias Matem\'aticas e de
  Computa\c{c}\~ao, Universidade de S\~ao Paulo},
  city={S\~ao Carlos}, state={SP}, country={Brazil}}

\begin{abstract}
\missing{abstract --- to be written last, once the eight runs of the matrix are
closed. What it has to say, if the measurements hold: that both methods were
built on the \emph{same} machinery --- same interpolated properties, same
momentum solver, coupling at the same point of the multiphase step --- so that
the measured difference is attributable to the interface representation and not
to the arrangement; that both reproduce the ellipsoidal regime of the Newtonian
reference case; that the two converge toward each other under refinement,
approaching from opposite sides; that at the reference resolution VOF conserves
mass about five times better, and that front-tracking's drift grows with
refinement; that the viscoelastic matrix reproduces the published orderings and
values, with VOF above the published numbers on all nine comparisons and
front-tracking at or below on seven; and what adaptive refinement costs and
returns in each.}
\end{abstract}

\begin{keyword}
front-tracking \sep volume of fluid \sep viscoelastic two-phase flow \sep
adaptive mesh refinement \sep hierarchical grids \sep PLIC
\end{keyword}

\end{frontmatter}

\tableofcontents

\section{Introduction}
\label{sec:intro}

\status{To be written. The bibliographic material is gathered --- the
comparison references and the parameters of every case are in
Sections~\ref{sec:newtonian} and \ref{sec:viscoelastic} --- but the framing
text has not been drafted.}

\missing{framing: the rising viscoelastic bubble as a problem; why the
interface representation matters when elastic stress concentrates in the wake;
the gap this work closes}

\subsection*{Contributions}

Four claims, in decreasing order of how firmly they are established:

\begin{enumerate}
\item \textbf{Three-dimensional front-tracking for viscoelastic two-phase
flow.} The reference implementation in the same solver~\cite{silva-thesis} uses
VOF. The triangulated surface with explicit connectivity, surface tension by a
per-triangle line integral, and mesh surgery is new, and it inherits the
viscoelastic path without reimplementing the constitutive equation.

\item \textbf{The comparison isolates the interface representation.} Both
methods share everything else --- $\rho(\mathbf{x})$, $\mu(\mathbf{x})$, the
variable-coefficient momentum solve, the projection --- and the front-tracking
coupling enters at the \emph{same point} of the multiphase step where VOF
advects the volume fraction. The measured difference is attributable to the
representation because nothing else differs. Section~\ref{sec:method} makes
that point precise.

\item \textbf{A full factorial design.} Two methods $\times$ two rheologies
$\times$ two meshes, all three-dimensional. The eight cells let the effect of
the interface representation be separated from the effect of the rheology and
from the effect of refinement, instead of being confounded with them.

\item \textbf{Adaptive refinement in both paths, with measured cost.}
\missing{depends on Section~\ref{sec:cost}; the 3D refinement criterion does not
exist yet}
\end{enumerate}

\noindent
One claim this work does \emph{not} make is that either method is more accurate
than the other. What the measurements support is \emph{where} each does better.

\section{Formulation}
\label{sec:formulation}

\status{Ready to write: the equations and the non-dimensionalisation convention
are established and checked against the code. The text below is complete on the
points that matter for reproducibility.}

\subsection{Governing equations}

Incompressible two-phase flow with surface tension and a polymeric
contribution. In non-dimensional form,
\begin{equation}
\rho\left(\frac{\partial \mathbf{u}}{\partial t}
          + \mathbf{u}\cdot\nabla\mathbf{u}\right)
= -\nabla p
  + \frac{1}{\mathrm{Re}}\nabla\cdot\!\left[\mu\left(\nabla\mathbf{u}
    + \nabla\mathbf{u}^{\mathsf{T}}\right)\right]
  + \rho\,\hat{\mathbf{g}}
  + \frac{1}{\mathrm{Bo}}\,\kappa\mathbf{n}\delta
  + \nabla\cdot\mathbf{S},
\label{eq:momentum}
\end{equation}
with
\begin{equation}
\mathrm{Re} = \frac{\rho_1 g^{1/2} L^{3/2}}{\mu_1},
\qquad
\mathrm{Bo} = \frac{\rho_1 g L^2}{\sigma},
\qquad
V = \sqrt{gL},
\qquad
L = 2r,
\label{eq:dimensionless}
\end{equation}
where $L$ is the bubble \emph{diameter} and the characteristic velocity is tied
to gravity. The volume fraction $f$ obeys
\begin{equation}
\frac{\partial f}{\partial t} + \nabla\cdot(\mathbf{u}f) = 0,
\label{eq:transport}
\end{equation}
and the properties are interpolated as $\rho = (1-f)\rho_0 + f\rho_1$ and
$\mu = (1-f)\mu_0 + f\mu_1$.

\subsection{The polymeric tensor is not the stored tensor}

This deserves its own subsection because it is a source of silent error. The
solver stores $\mathbf{S}$; the polymeric tensor is
\begin{equation}
\mathbf{T} = \mathbf{S} + \frac{2(1-\beta)}{\mathrm{Re}}\,\mathbf{D},
\qquad
\mathbf{D} = \tfrac{1}{2}\!\left(\nabla\mathbf{u}
           + \nabla\mathbf{u}^{\mathsf{T}}\right),
\label{eq:T-from-S}
\end{equation}
and $\beta$ is the \emph{interpolated} value between the phases, by the same
expression the solver uses:
\begin{equation}
\beta = \frac{(1-f)\,\mu_0\beta_0 + f\,\mu_1\beta_1}{\mu}.
\label{eq:beta-interp}
\end{equation}
Mistaking $\mathbf{S}$ for $\mathbf{T}$ raises no error; it produces a
plausible and wrong field. The output files of this work carry both, labelled.

\subsection{Constitutive model}

Oldroyd-B, as an evolution equation for the conformation tensor $\mathbf{A}$:
\begin{equation}
\mathrm{De}\left(\frac{\partial \mathbf{A}}{\partial t}
  + \mathbf{u}\cdot\nabla\mathbf{A}
  - \nabla\mathbf{u}\cdot\mathbf{A}
  - \mathbf{A}\cdot\nabla\mathbf{u}^{\mathsf{T}}\right)
= \mathbf{I} - \mathbf{A},
\qquad
\mathrm{De} = \frac{U\lambda}{L}.
\label{eq:oldroyd}
\end{equation}
Oldroyd-B has no finite extensibility limit, so the stress grows without bound
under sustained extension. That is a property of the model, not of the
discretisation, and it bounds how far the viscoelastic cases can be pushed ---
Section~\ref{sec:viscoelastic} returns to it.

\subsection{The solver's non-dimensionalisation convention}
\label{sec:convention}

This subsection exists for reproducibility, and its content was \emph{read from
the code and checked numerically}, not assumed. The solver integrates the
equations in \textbf{dimensional} variables, with the dimensionless groups set
so that the coefficients of (\ref{eq:momentum}) reduce to the physical values:
\begin{equation}
\mathrm{Re} = 1,
\qquad
\mathrm{Fr} = 1/\sqrt{g},
\qquad
\mathrm{Ca} = 1/\sigma,
\end{equation}
with $\rho$ and $\mu$ entering at their tabulated values. The gravitational
term appears as $1/\mathrm{Fr}^2$ and the interfacial one as
$1/(\mathrm{Ca}\cdot\mathrm{Re})$. Checked: $\mathrm{Fr} = 1.010153$ gives
$1/\mathrm{Fr}^2 = 0.979999$, which is $g = 0.98$; and
$\mathrm{Ca} = 0.0408163265$ gives $1/(\mathrm{Ca}\cdot\mathrm{Re}) = 24.5$,
which is $\sigma$. With $U = L = 1$ at this scaling, $\mathrm{De} = \lambda$.

\section{Numerical method}
\label{sec:method}

\status{Ready to write: every part described here is implemented and verified.
The subsection on adaptive refinement (\ref{sec:amr}) is the exception --- the
criterion exists and is measured in two dimensions, but has not been ported to
three.}

\subsection{Shared machinery, and the single point of difference}
\label{sec:shared}

This subsection carries the methodological claim of the paper, so it states the
arrangement precisely.

Both methods use the same projection method, the same variable-coefficient
momentum solve, and the same property interpolation. The front-tracking
coupling is a function pointer, null by default, invoked inside the multiphase
step \emph{after} the source terms and \emph{before} the viscosity and density
interpolation. At that point it overwrites the volume fraction from the
geometry of the front, so the properties are built from the new fraction. VOF
reaches the same field at the same point by advecting (\ref{eq:transport}) with
a geometric PLIC reconstruction.

Everything downstream of that point --- $\rho(\mathbf{x})$, $\mu(\mathbf{x})$,
the momentum solve, the pressure projection, the output probes --- is
bit-identical code. The consequence is that a measured difference between the
two is a difference in how the interface is represented, and not in the solver,
the mesh, the time step, the initial geometry, or the instrument.

\subsection{VOF: geometric PLIC}

\missing{PLIC reconstruction and the split advection --- short, since it is the
pre-existing method; cite the reference implementation}

\paragraph{Initial volume fraction.}
The initial fraction is computed by \textbf{sub-cell subdivision}: each
interface cell is split into $16^3$ sub-cells, each resolved by the closed-form
plane--box intersection volume, and the sum is condensed onto the coarse cell.
This matters more than it appears. A one-cell smoothed step --- the obvious
shortcut --- gives a bubble volume error of \num{1.23e-2} at the running
resolution against \num{6.11e-5} for the subdivision, and its measured order of
convergence is erratic (1.59, 2.90, 0.99) because the ramp width is tied to $h$.
Since VOF conserves faithfully whatever mass it is given, an initial-condition
error of that size never recovers.

\subsection{Front-tracking: explicit triangulated surface}

The interface is a closed triangulated surface with explicit connectivity,
built as an icosphere by recursive subdivision of an icosahedron.

\paragraph{Surface tension by line integral.}
The force on each triangle is obtained from the line integral along its edges,
$\mathbf{f} = \sigma(\mathbf{b}-\mathbf{a})\times\mathbf{n}$, distributed to the
two endpoints. The force is $-\sigma\nabla A$ --- tension minimises area ---
and the sign was derived rather than guessed, then checked against the numerical
area gradient.

\paragraph{The Laplace oracle in three dimensions.}
For a sphere the pressure jump is $\Delta p = 2\sigma/R$, \emph{not} $\sigma/R$.
The two-dimensional value is $\sigma/R$, and a 2D implementation that closes to
0.02\% against it would, if the number were carried over, validate the wrong
formula. The distinction is flagged in the header, the harness and the example.

\paragraph{Surgery.}
Split, collapse and flip, with a curved midpoint
\begin{equation}
\mathbf{m} = \frac{\mathbf{a}+\mathbf{b}}{2}
  - \frac{1}{8}\Big[\big((\mathbf{b}-\mathbf{a})\cdot\mathbf{n}_a\big)\mathbf{n}_a
  + \big((\mathbf{a}-\mathbf{b})\cdot\mathbf{n}_b\big)\mathbf{n}_b\Big],
\end{equation}
which is the PN-triangle (Hermite) construction. Surgery keeps the triangle
edge matched to the cell size as the bubble deforms.

\paragraph{Volume fraction from the surface.}
The fraction is obtained from the signed distance and normal to the surface,
then from the closed-form plane--box volume. The sign comes from parity
ray-casting, never from the nearest triangle's normal --- the latter is wrong
for a cell straddling a concavity.

\subsection{Interface extraction for comparison}
\label{sec:mc}

VOF has no interface as an object; the comparison needs one. The isosurface
$f = 1/2$ is extracted by marching cubes on the \emph{dual} grid --- the volume
fraction lives at cell centres, so the lattice nodes are cell centres --- with
vertices \textbf{welded} by lattice edge, which makes the extracted surface
continuous. The alternative, writing the per-cell PLIC planes, gives a surface
that is discontinuous from one cell to the next; that discontinuity is the right
object for auditing the reconstruction and the wrong one for measuring shape.

Both methods therefore emit a triangulated surface in the same file format, and
every shape quantity in this paper is computed from those surfaces by the
\emph{same} routine. The instrument is common; only the surface differs.

\paragraph{On adaptive meshes the extraction runs on the finest level only.}
This is sufficient, and not by economy: the band criterion refines to the finest
level with a floor of five fine cells on each side of the interface, and outside
that band the fraction is exactly 0 or 1. The eight-node stencil therefore never
reaches the seam between levels, which is what normally breaks isosurface
extraction on adaptive meshes.

\subsection{Sub-grid undulation, and the filter}
\label{sec:tsur}

\status{Implemented, verified, and \emph{measured on the production runs}. The
with/without comparison at the reference resolution is pending.}

A Lagrangian surface accumulates error as vertex-scale zig-zag, because the
vertices move independently under the interpolated velocity and nothing
constrains the sub-grid shape. Quantifying it needs a shape-independent
indicator, and the natural one is already computed by the filter itself: for an
edge $(v_1,v_2)$ with edge normal $\mathbf{n}_e$ and midpoint $\mathbf{m}$,
\begin{equation}
h_1 = (\mathbf{v}_1 - \mathbf{m})\cdot\mathbf{n}_e ,
\end{equation}
and because $\mathbf{m}$ is the midpoint, $h_2 = -h_1$ exactly --- so $h_1$ is
the edge's \emph{tilt} relative to the local surface. Normalised by the
half-edge it is dimensionless, and on a smooth surface it is residual whatever
the shape. Measuring radial deviation instead, as the synthetic tests do, would
measure deformation on a deformed bubble.

\begin{table}[htbp]
\centering
\caption{Undulation at the final frame. ``Above'' is the fraction of edges with
indicator above $0.05$; the initial sphere is the oracle.}
\label{tab:undulation}
\small
\begin{adjustbox}{max width=\linewidth}
\pgfplotstabletypeset[
  col sep=space, string type,
  columns/caso/.style={column name={Run}, column type={l}},
  columns/acima/.style={column name={above}, column type={r}},
  columns/inferior/.style={column name={lower half}, column type={r}},
  columns/mediana/.style={column name={median}, column type={r}},
  columns/maximo/.style={column name={max}, column type={r}},
  every head row/.style={before row=\toprule, after row=\midrule},
  every last row/.style={after row=\bottomrule},
]{\dat{ondulacao.dat}}
\end{adjustbox}
\end{table}

Two things in that table. An indicator of $0.997$ means an edge nearly
\emph{perpendicular} to the local surface --- a full zig-zag. And the undulation
is not diffuse: of the 1289 edges above threshold in the coarse Newtonian run,
\textbf{all 1289} lie in the lower half of the bubble. In the worst case, 1765
of 1770.

\subsubsection{The two methods differ in kind, not degree}

\begin{table}[htbp]
\centering
\caption{The indicator over time, same case, both methods. The VOF
\emph{level} is not comparable --- see below.}
\label{tab:undtime}
\small
\pgfplotstabletypeset[
  col sep=space, string type,
  columns/t/.style={column name={$t$}, column type={r}},
  columns/ft/.style={column name={FT above (\%)}, column type={r}},
  columns/ftmax/.style={column name={FT max}, column type={r}},
  columns/vof/.style={column name={VOF above (\%)}, column type={r}},
  columns/vofmax/.style={column name={VOF max}, column type={r}},
  every head row/.style={before row=\toprule, after row=\midrule},
  every last row/.style={after row=\bottomrule},
]{\dat{ondulacao-tempo.dat}}
\end{table}

The VOF \emph{absolute} level means nothing: it is already 27\% at $t=0$, where
the surface is the same sphere front-tracking represents with $0.00\%$. That is
faceting of the marching-cubes extraction, not of the interface. Only the trend
is comparable, and the trends are opposite --- front-tracking accumulates from
zero to 21\% with the maximum reaching $0.997$; VOF stays at $32\pm2\%$ with the
maximum never above $0.37$.

The reason is structural. PLIC reconstructs \emph{one plane per cell}, and a
plane cannot undulate within a cell; the reconstruction is re-derived from the
field every step, so nothing accumulates at the surface level. Put the other
way: \textbf{the PLIC reconstruction is itself a filter}, and a far more
aggressive one than the method below --- it re-imposes a planar interface every
step. The price is the dual of the benefit: it cannot represent anything below
the cell, which is exactly what a cusp is.

So the two methods sit at opposite ends of one trade-off: VOF cannot undulate
and cannot resolve sub-grid curvature; front-tracking can resolve it and
contaminates it with undulation at the same scale. That makes the filter the
difference between the two methods being compared on the cusp, rather than one
of them being disqualified from it.

\subsubsection{The filter, and why it needs a criterion}

BALANCE repositions each vertex along its ring normal to the height that
\emph{reproduces} the star's volume; TSUR puts both ends of an edge at a common
height, weighted by the two star volumes. Both are from the reference
Freeflow implementation~\cite{desousa2004}, where the procedure is called a
``mass-conserving filter''.

That property is measured, not taken on faith: on a sphere with an $8\%$ radial
undulation, three passes drop the deviation by $84\%$ with a volume change of
$0.000\%$, where a Laplacian average smooths slightly \emph{more} and loses
$9.74\%$ of the volume.

\paragraph{But conserving volume does not protect against artificial surface
tension.} Reducing area at fixed volume is precisely what surface tension does.
On a \emph{sphere} the effect is invisible, because the smooth state is already
the minimum-area shape and the filter pushes where the physics also pushes. An
ellipsoid exposes it: at fixed volume, 20 passes take the axis ratio from
$2.0239$ to $1.9072$ and the sphericity from $0.925670$ to $0.931380$,
monotonically, with $\mathrm{d}V/V = 2\times10^{-15}$.

That is $0.00029$ of sphericity per pass. Applied every 20 steps over 1501
steps --- 75 applications --- it would induce $0.021$, which is \textbf{nine
times} the sphericity difference between the two methods that this paper
measures ($0.0023$). Unconditional application would destroy the measurement.

\paragraph{The criterion.} Apply only where there \emph{is} undulation, gated on
the indicator above. There is no trade-off to tune; the plateau is wide.

\begin{table}[htbp]
\centering
\caption{Threshold sweep. The second column must stay far below $0.0023$; the
last must stay high.}
\label{tab:threshold}
\small
\adjustbox{max width=\linewidth}{%
\pgfplotstabletypeset[
  col sep=space, string type,
  columns/limiar/.style={column name={threshold}, column type={r}},
  columns/dpsi/.style={column name={$\Delta\Psi$, 75 passes, smooth ellipsoid}, column type={r}},
  columns/movidos/.style={column name={vertices moved/pass}, column type={r}},
  columns/queda/.style={column name={undulation removed}, column type={r}},
  every head row/.style={before row=\toprule, after row=\midrule},
  every last row/.style={after row=\bottomrule},
]{\dat{filtro-limiar.dat}}
}
\end{table}

At any threshold above $0.02$ the smooth ellipsoid is left \emph{completely
untouched} --- zero vertices moved per pass --- and the undulation still drops
by $82\%$. We use $0.05$, with margin from the edge.

\missing{the production with/without pair at the reference resolution. The
without-filter half is measured (Table~\ref{tab:undtime}); the with-filter run
is pending.}

\subsection{Adaptive refinement}
\label{sec:amr}

\status{Implemented and verified in three dimensions, including the dynamic
remeshing loop. No production run on an adapted mesh yet.}

The criterion is the maximum of two: a geometric band around the interface,
derived from a minimum-fine-cell rule, and a vorticity threshold that captures
the wake. Leaving the fine mesh is deliberately harder than entering it --- a
cell already refined only coarsens once $|\omega|$ falls below \emph{half} the
threshold. Without that hysteresis the two-dimensional work measured flapping:
cells crossing a bare threshold every cycle and the mesh drifting in size.

Three things differ from the two-dimensional case, and the first is not a
detail.

\paragraph{The level table is built by dilation, not by search.}
In two dimensions the distance from every cell to the nearest seed is found by
brute force, $O(n_{\text{cell}} \times n_{\text{seed}})$, and that is affordable.
In three dimensions, 250\,000 cells against thousands of seeds is of order
$10^9$ operations \emph{per remesh}. Instead each seed stamps its neighbourhood
of radius $\lceil \mathrm{thr}_0/h \rceil$, which is
$O(n_{\text{seed}} \times R^3)$.

The dilation is \emph{exact}, not an approximation: the level is monotone
decreasing in distance, so the maximum over seeds is the level of the nearest
seed. That argument is not taken on trust --- the exhaustive version is kept
alongside and the harness requires byte-for-byte equality in five
configurations. Zero cells differ in any of them.

\paragraph{Vorticity stops being a scalar.}
$\omega_z$ suffices in the plane; in space the criterion uses
$|\nabla \times \mathbf{u}|$ with all three components, which costs twelve facet
evaluations per cell against four.

\paragraph{The mesh file format.}
Generic in the dimension, but with six fields per patch instead of four, and
one-based indices. In two dimensions the first version used zero-based indices:
the count of fine cells came out \emph{right} and the positions \emph{wrong},
displaced by one cell of the level, and the measured band fell from 5 to 1.5
cells. A right count with wrong positions appears in no total, so the harness
re-reads the file and checks every patch against the table.

\paragraph{What the loop delivers.}
Measured, base 25 with one level, remeshing every ten steps: the band on the
mesh the solver actually built is 6.03 fine cells against a floor of five; no
kernel support point lands in a cell of different size; no position is lost in
the transfer. The last two must be exactly zero --- the kernel is normalised for
a single $h$, and crossing a level breaks the partition of unity \emph{silently},
because the weights can still sum to nearly one.

\section{Verification}
\label{sec:verification}

\status{Ready. Every entry in Table~\ref{tab:oracles} is measured in this tree.
What is \emph{not} covered is stated at the end of the section, and that
statement is part of the result.}

The verification is organised around a distinction that the rest of the paper
relies on: an oracle is either \textbf{exact}, in which case it passes or fails
with no tolerance, or \textbf{convergent}, in which case what is checked is the
measured order in $h$. Mixing the two is how a plausible-but-wrong
implementation survives a test suite --- a tolerance chosen after seeing the
answer is not a test.

\begin{table}[htbp]
\centering
\caption{Verification oracles. Exact oracles admit no tolerance; convergent
ones are checked by measured order.}
\label{tab:oracles}
\small
\begin{adjustbox}{max width=\linewidth}
\pgfplotstabletypeset[
  col sep=space, string type,
  columns={peca,oraculo,resultado},
  columns/peca/.style={column name={Component}, column type={l}},
  columns/oraculo/.style={column name={Oracle}, column type={l}},
  columns/resultado/.style={column name={Result}, column type={l}},
  every head row/.style={before row=\toprule, after row=\midrule},
  every last row/.style={after row=\bottomrule},
]{\dat{oraculos.dat}}
\end{adjustbox}
\end{table}

\subsection{Two oracles worth describing}

\paragraph{The marching-cubes table is checked against geometry, not against
another table.}
The 256-case table is copied data, and a single wrong digit produces a hole in
the surface that passes easily as a rendering artefact. Rather than compare
against a second hand-written table --- two copies of the same error would agree
--- the harness \emph{derives}, from each case, which edges must be cut (their
endpoints differ in sign) and requires the set the table uses to be exactly that
set. Zero of 256 cases are inconsistent.

\paragraph{The torus is in the suite on purpose.}
A sphere exercises few of the 256 cases. A torus exercises many more, and
$\chi = 0$ is a sharp number: a single missing triangle moves it. Two separated
spheres give $\chi = 4$ for the same reason.

\subsection{What the green does not prove}

Three gaps, stated because a verification section that lists only its successes
is advertising:

\begin{itemize}
\item \textbf{No oracle covers more than one process.} The parallel path has
never been measured in three dimensions, and the interface extraction is serial
by construction.
\item \textbf{No oracle covers adaptive refinement in three dimensions},
because it does not exist yet (Section~\ref{sec:amr}).
\item \textbf{The comparison against the literature is against published
figures, not tables.} Agreement on shape is therefore qualitative; only the
centre of mass and the rise velocity are curves we can read.
\end{itemize}

\section{Newtonian three-dimensional bubble}
\label{sec:newtonian}

\status{Uniform mesh is \textbf{closed at both resolutions} --- $25\times25\times50$
and the reference $50\times50\times100$, 1501 steps each, in both methods. Case~2
has not been run. The adaptive half now has a working dynamic remeshing loop
(Section~\ref{sec:amr}) but no production run.}

\subsection{Case definition}

A bubble of radius $R = 0.25$ centred at $(0.5,0.5,0.5)$ in a tank
$[0,1]\times[0,1]\times[0,2]$, no-slip on all walls, zero initial velocity,
pressure and polymeric tensor, integrated to $t = 3$. Two parameter sets,
Table~\ref{tab:params}, taken from the reference work and originally due to
Adelsberger et al.

\begin{table}[htbp]
\centering
\caption{Physical and dimensionless parameters.}
\label{tab:params}
\small
\adjustbox{max width=\linewidth}{%
\pgfplotstabletypeset[
  col sep=space, string type,
  every head row/.style={before row=\toprule, after row=\midrule},
  every last row/.style={after row=\bottomrule},
]{\dat{caso-parametros.dat}}
}
\end{table}

\subsection{Uniform mesh}

\subsubsection{The methods converge toward each other}

This is the central result of the section, and it is a statement about both
methods rather than about either one: refining the mesh brings them together.

\begin{table}[htbp]
\centering
\caption{Difference between the two methods at $t=3$, at both resolutions.}
\label{tab:convergence}
\small
\begin{adjustbox}{max width=\linewidth}
\pgfplotstabletypeset[
  col sep=space, string type,
  columns/grandeza/.style={column name={Quantity}, column type={l}},
  columns/n25/.style={column name={$25\times25\times50$}, column type={r}},
  columns/n50/.style={column name={$50\times50\times100$}, column type={r}},
  columns/fator/.style={column name={Factor}, column type={r}},
  every head row/.style={before row=\toprule, after row=\midrule},
  every last row/.style={after row=\bottomrule},
]{\dat{c1-convergencia.dat}}
\end{adjustbox}
\end{table}

What makes this more than a coincidence is the \emph{direction} each method
moves. The VOF centre of mass falls from $1.4903$ to $1.4766$; the
front-tracking one rises from $1.4615$ to $1.4678$. They approach from opposite
sides toward roughly $1.47$. A single method converging says only that it is
consistent; two independent representations converging from opposite sides
toward the same value is evidence that the value is the right one.

\begin{table}[htbp]
\centering
\caption{Front-tracking against VOF at $t=3$, both resolutions. Shape
quantities from the extracted surfaces by the same routine; mass from the field
integral for VOF and from the surface volume for front-tracking.}
\label{tab:resolution}
\small
\begin{adjustbox}{max width=\linewidth}
\pgfplotstabletypeset[
  col sep=space, string type,
  columns/grandeza/.style={column name={Quantity at $t=3$}, column type={l}},
  columns/vof25/.style={column name={VOF $N{=}25$}, column type={r}},
  columns/ft25/.style={column name={FT $N{=}25$}, column type={r}},
  columns/vof50/.style={column name={VOF $N{=}50$}, column type={r}},
  columns/ft50/.style={column name={FT $N{=}50$}, column type={r}},
  every head row/.style={before row=\toprule, after row=\midrule},
  every last row/.style={after row=\bottomrule},
]{\dat{c1-resolucao.dat}}
\end{adjustbox}
\end{table}

Both methods reproduce the \textbf{ellipsoidal regime} the reference reports for
case~1 --- the height-to-equivalent-diameter ratio is below one in both, that
is, flattened along the direction of rise --- and both reproduce the shape of
the published rise-velocity curve: the velocity grows, peaks near $t = 0.9$, and
decreases toward a steady value.

\begin{figure}[htbp]
\centering
\begin{tikzpicture}
\begin{axis}[width=0.47\linewidth, xlabel={$t$}, ylabel={$z_c$},
             legend pos=north west, grid=both, no markers, legend cell align=left]
  \addplot table[x=t, y=zc] {\dat{bolha3d-c1-N25-vof.dat}}; \addlegendentry{VOF $N{=}25$}
  \addplot table[x=t, y=zc] {\dat{bolha3d-c1-N25-ft.dat}};  \addlegendentry{FT $N{=}25$}
  \addplot table[x=t, y=zc] {\dat{bolha3d-c1-N50-vof.dat}}; \addlegendentry{VOF $N{=}50$}
  \addplot table[x=t, y=zc] {\dat{bolha3d-c1-N50-ft.dat}};  \addlegendentry{FT $N{=}50$}
\end{axis}
\end{tikzpicture}\hfill
\begin{tikzpicture}
\begin{axis}[width=0.47\linewidth, xlabel={$t$}, ylabel={$\Psi$},
             legend pos=south west, grid=both, no markers, legend cell align=left]
  \addplot table[x=t, y=Psi] {\dat{bolha3d-c1-N25-vof.dat}}; \addlegendentry{VOF $N{=}25$}
  \addplot table[x=t, y=Psi] {\dat{bolha3d-c1-N25-ft.dat}};  \addlegendentry{FT $N{=}25$}
  \addplot table[x=t, y=Psi] {\dat{bolha3d-c1-N50-vof.dat}}; \addlegendentry{VOF $N{=}50$}
  \addplot table[x=t, y=Psi] {\dat{bolha3d-c1-N50-ft.dat}};  \addlegendentry{FT $N{=}50$}
\end{axis}
\end{tikzpicture}
\caption{Centre of mass and sphericity at both resolutions.}
\label{fig:resolution}
\end{figure}

\subsubsection{Mass conservation goes the other way}

Mass is the one quantity that does \emph{not} improve for both methods, and the
two resolutions disagree about which method is better --- which is itself the
finding.

At $25\times25\times50$ the front-tracking drift is $-0.096\%$ against
$+0.396\%$ for VOF, a factor of four in favour of front-tracking. At
$50\times50\times100$ it reverses: $-0.851\%$ against $-0.170\%$, a factor of
five the other way.

Reading the time history explains it. Up to step 400 the two front-tracking runs
are nearly identical ($-7.3\times10^{-4}$ and $-6.8\times10^{-4}$ at step 200,
$-2.5\times10^{-3}$ and $-2.8\times10^{-3}$ at step 400). After that they
diverge: the coarse run peaks at $-2.6\times10^{-3}$ and \emph{recovers} to
$-9.6\times10^{-4}$, while the fine run accumulates monotonically. The VOF drift
likewise changes sign between the resolutions. Both coarse-mesh figures are
therefore partly \textbf{cancellation}, not conservation, and neither should be
quoted as a property of its method.

What survives is the fine-mesh statement: at the reference resolution VOF
conserves mass about five times better than front-tracking, and front-tracking's
drift \emph{grows} with refinement.

\subsubsection{The drift that grows is an artefact of the comparison}
\label{sec:deriva-nsub}

The observation that the triangulation ends with the same vertex and triangle
counts at both resolutions is not a puzzle accompanying the drift: it
\emph{is} the drift's explanation. Both production runs were launched with the
same surface subdivision, so the mesh was refined by a factor of two and the
front was not. The ratio $ds/h$ therefore \emph{doubled} between the two runs,
and the pair never was a mesh-refinement comparison --- it varied two things at
once, and the one it was meant to isolate was not the one that moved.

A controlled pair settles it. Running the same case on two meshes with the
surface resolution \emph{following} the grid, so that $ds/h = 0.75$ in both
($N=20$ with three subdivisions, $N=40$ with four), the drift falls:

\adjustbox{max width=\linewidth}{%
\pgfplotstabletypeset[
  string type,
  columns/passo/.style={column name={step}},
  columns/aa0/.style={column name={$A/A_0$}},
  columns/n20/.style={column name={$N=20$}},
  columns/cir20/.style={column name={surg.}},
  columns/n40/.style={column name={$N=40$}},
  columns/cir40/.style={column name={surg.}},
  columns/razao/.style={column name={ratio}},
  every head row/.style={before row=\toprule, after row=\midrule},
  every last row/.style={after row=\bottomrule},
]{\dat{ft-duas-resolucoes.dat}}
}

The ratio stays between $3.2$ and $3.7$ throughout; second order would give
$4$, first order $2$. Over the three points before the first surgery --- pure
advection against pure advection --- it is $3.63$, $3.70$ and $3.61$, a mean of
$3.65$, or an observed order of $\mathbf{1.87}$. \textbf{The drift is
discretisation error and vanishes under refinement.} There is no systematic
component: an interpolation bias would not fall with $h$.

\caveat{At step 500 the ratio drops to $3.21$, and there the fine case has
performed 252 surgery operations against the coarse case's 55 --- the
comparison is no longer advection against advection alone. The clean points are
the early ones, and they are the ones that carry the conclusion.}

Surgery was the named suspect, and a second experiment with everything else
held identical acquits it:

\adjustbox{max width=\linewidth}{%
\pgfplotstabletypeset[
  string type,
  columns/passo/.style={column name={step}},
  columns/comcirurgia/.style={column name={with surgery}},
  columns/semcirurgia/.style={column name={without surgery}},
  columns/razao/.style={column name={worse by}},
  every head row/.style={before row=\toprule, after row=\midrule},
  every last row/.style={after row=\bottomrule},
]{\dat{ft-cirurgia.dat}}
}

\textbf{Surgery does not cause the drift; it reduces it by $42\%$.} Without it
the triangles stretch and the deformed surface is worse represented. The two
runs are \emph{identical up to step 200}, which is where the first surgery
occurs --- the check that the only difference between them is the intended one.

What this leaves standing is narrower and more useful than the original claim:
at the reference resolution VOF conserves mass about five times better, and
front-tracking's drift is a discretisation error of order $1.87$ that the
production pair failed to expose because it held the surface resolution fixed
while refining the mesh.

\subsubsection{The shapes}

The quantities above summarise the shape; the shape itself is a planar section.
Figure~\ref{fig:newtshape} cuts each extracted surface by the vertical plane
through the bubble centre --- the same view the reference uses to compare two
codes --- and overlays the two methods.

The contours are drawn \emph{centred on their own centroid}. Overlaying them at
their absolute heights would hide the difference in shape behind the difference
in rise, and the rise is already in Table~\ref{tab:resolution}.

\begin{figure}[htbp]
\centering
\begin{tikzpicture}
\begin{axis}[width=0.46\linewidth, height=0.42\linewidth,
             xlabel={$x - x_c$}, ylabel={$z - z_c$}, axis equal image,
             title={$25\times25\times50$}, grid=both, no markers,
             legend pos=south east, legend cell align=left, font=\small]
  \addplot+[thick] table[x=x, y=z] {\dat{secoes/newt-N25-vof-c.dat}};
    \addlegendentry{VOF}
  \addplot+[thick, dashed] table[x=x, y=z] {\dat{secoes/newt-N25-ft-c.dat}};
    \addlegendentry{FT}
\end{axis}
\end{tikzpicture}\hfill
\begin{tikzpicture}
\begin{axis}[width=0.46\linewidth, height=0.42\linewidth,
             xlabel={$x - x_c$}, ylabel={$z - z_c$}, axis equal image,
             title={$50\times50\times100$}, grid=both, no markers,
             legend pos=south east, legend cell align=left, font=\small]
  \addplot+[thick] table[x=x, y=z] {\dat{secoes/newt-N50-vof-c.dat}};
    \addlegendentry{VOF}
  \addplot+[thick, dashed] table[x=x, y=z] {\dat{secoes/newt-N50-ft-c.dat}};
    \addlegendentry{FT}
\end{axis}
\end{tikzpicture}
\caption{Newtonian case 1 at $t=3$: vertical section through the bubble centre,
front-tracking against VOF, at both resolutions. Contours centred on their own
centroid so that the comparison is of shape.}
\label{fig:newtshape}
\end{figure}

Both are oblate, as the ellipsoidal regime requires. The visible difference is
not in the overall form but in the \emph{resolution of the contour}: the
marching-cubes section carries 97 points at the coarse mesh against 185 for the
triangulated surface, because the extraction lattice is the fluid mesh while the
triangulation is refined by surgery independently of it. At
$50\times50\times100$ the two carry 188 and 196, and the contours lie close
enough that the faceting of each is the larger visible effect.

\paragraph{The instrument was cleared, not assumed.}
The VOF centre of mass is computed from the marching-cubes surface, which
encloses about 2.3\% less volume than the field at the coarse resolution; that
could bias the centroid. An independent probe computes the centroid directly
from the field and gives $1.490212$ against $1.4903$ from the surface ---
agreement to \num{1e-4}. The volume bias does not translate into a position
bias.

\paragraph{Topological soundness.}
121 output frames per method per resolution, \emph{all} surfaces closed (zero
orphan edges, $\chi = 2$). In the front-tracking runs the triangulation grew from
5120 to 5558 triangles, with 8824 and 9604 surgery operations and a
partition-of-unity residual of \num{2.1e-14}.

\subsection{Case 2: topological change}
\missing{$\rho_1/\rho_2 = 1000$, $\mu_1/\mu_2 = 100$, $\sigma = 1.96$. Lower
surface tension favours break-up, which is where an explicit surface with
surgery and a volume-fraction field should differ most. Not run.}

\subsection{Adaptive mesh}
\missing{blocked on Section~\ref{sec:amr}}

\section{Viscoelastic three-dimensional bubble}
\label{sec:viscoelastic}

\status{\textbf{Closed on the uniform mesh}: all twelve runs of the matrix
(six $(\lambda_1,\beta)$ combinations $\times$ two methods) completed 2001 steps
at $25\times25\times75$ with zero non-finite values. The reference resolution is
complete for the cusp case in VOF and running in front-tracking. The adaptive
half has a working remeshing loop but no production run.

One claim in this section was \emph{withdrawn} by its own measurement: the
velocity drop is not systematically larger in VOF. See
Section~\ref{sec:dropdef}.}

\subsection{Case definition}

The Newtonian case of Section~\ref{sec:newtonian} with an Oldroyd-B medium
surrounding a Newtonian bubble, following the reference
thesis~\cite{silva-thesis}. Two changes to the geometry and the schedule, both
with a reason:

\begin{itemize}
\item the tank is raised to $[0,1]\times[0,1]\times[0,3]$, so the top and bottom
walls do not influence the polymeric tensor;
\item the final time goes from $t=3$ to $t=4$, because the elastic stress takes
longer than the velocity to reach a steady state.
\end{itemize}

\begin{table}[htbp]
\centering
\caption{Viscoelastic matrix: relaxation time and solvent fraction.}
\label{tab:vematrix}
\small
\begin{adjustbox}{max width=\linewidth}
\pgfplotstabletypeset[
  col sep=space, string type,
  columns/lambda/.style={column name={$\lambda_1$}, column type={c}},
  columns/beta/.style={column name={$\beta$}, column type={c}},
  columns/polimero/.style={column name={Polymer}, column type={c}},
  columns/esperado/.style={column name={Expected shape at $t=4$}, column type={l}},
  every head row/.style={before row=\toprule, after row=\midrule},
  every last row/.style={after row=\bottomrule},
]{\dat{ve-matriz.dat}}
\end{adjustbox}
\end{table}

\subsection{The cusp as the discriminating case}

At $\lambda_1 = 0.4$ and $\beta = 0.2$ the reference reports a \textbf{cusp}
forming at the trailing edge of the bubble: the rising motion deforms the
viscoelastic medium, which exerts elastic force back on the interface, and that
force stretches the rear of the bubble into a point.

This is the case that gives the comparison content rather than making it a tie.
A cusp is a region of high curvature concentrated on a scale that approaches the
cell size, and the two representations should handle it differently: an explicit
surface with surgery can refine locally along the interface, whereas a volume
fraction with a per-cell planar reconstruction cannot represent curvature below
one cell.

\status{\textbf{The pair is closed at the reference resolution.} Both methods
completed 2001 steps to $t=4$ with 161 extracted surfaces, all closed. The
front-tracking run carries the filter of Section~\ref{sec:tsur}, which is what
makes it comparable on this case at all.}

\begin{table}[htbp]
\centering
\caption{The pair at $t=4$, $50{\times}50{\times}150$, same instrument for both.}
\label{tab:ftn50}
\small
\begin{adjustbox}{max width=\linewidth}
\pgfplotstabletypeset[
  col sep=space, string type,
  columns/grandeza/.style={column name={}, column type={l}},
  columns/ft/.style={column name={front-tracking}, column type={r}},
  columns/vof/.style={column name={VOF}, column type={r}},
  every head row/.style={before row=\toprule, after row=\midrule},
  every last row/.style={after row=\bottomrule},
]{\dat{ft-n50-final.dat}}
\end{adjustbox}
\end{table}

\subsubsection{Two measures, and only one of them is free of the mesh}

The cusp is read with two numbers, and the pair exists because they fail
differently.

The \textbf{asymmetry} $A = (z_c - z_{\min})/(z_{\max} - z_c)$ is $1$ for a
sphere and grows as the rear is drawn out. It is \emph{integrated}, so no cell
size bounds it.

The \textbf{radius of curvature} $R$ at the lowest vertex --- a least-squares
sphere through that vertex and its ring --- is \emph{local}, and a method that
reconstructs one plane per cell cannot report a radius below the cell. It
therefore has a floor, and the floor is the point.

\begin{table}[htbp]
\centering
\caption{The two measures over time, both methods, reference resolution. The
asymmetry agrees throughout; $R/R_{\mathrm{eq}}$ does not, until both reach the
mesh floor.}
\label{tab:cusppair}
\small
\begin{adjustbox}{max width=\linewidth}
\pgfplotstabletypeset[
  col sep=space, string type,
  columns={t,ftA,vofA,dA,ftR,vofR},
  columns/t/.style={column name={$t$}, column type={r}},
  columns/ftA/.style={column name={FT $A$}, column type={r}},
  columns/vofA/.style={column name={VOF $A$}, column type={r}},
  columns/dA/.style={column name={diff.}, column type={r}},
  columns/ftR/.style={column name={FT $R/R_{\mathrm{eq}}$}, column type={r}},
  columns/vofR/.style={column name={VOF $R/R_{\mathrm{eq}}$}, column type={r}},
  every head row/.style={before row=\toprule, after row=\midrule},
  every last row/.style={after row=\bottomrule},
]{\dat{cuspide-par.dat}}
\end{adjustbox}
\end{table}

\paragraph{The asymmetry is the result.} The two methods track each other to
$0.2$--$2.4\%$ from $t=1.6$ to $t=4$, ending at $1.2846$ and $1.2907$ --- a
difference of $0.5\%$. On the case the reference singles out as discriminating,
with the two interface representations at their most different, the integrated
shape of the cusp is the same.

\paragraph{But their agreement in $R$ at $t=4$ is not convergence.} At $t=4$
front-tracking gives $R = 0.02142$ and VOF $R = 0.02168$, within $1\%$ --- and
$h = 1/50 = 0.02$, so \textbf{both are at $R \approx 1.07\,h$: one cell}. Two
methods pinned at the same floor agree for a reason that has nothing to do with
the physics, and reading that $1\%$ as convergence would be reading the mesh.

The informative part of that column is the \emph{approach}, not the endpoint.
VOF drops to $\approx 0.11$ by $t=2.4$ and then sits there for the rest of the
run --- it reaches its floor early and stays. Front-tracking descends
gradually, $0.40 \to 0.22 \to 0.13 \to 0.086$, arriving at the floor only at
$t=4$. That is the expected signature of the structural difference argued in
Section~\ref{sec:tsur}: a per-cell planar reconstruction saturates at the cell,
while an explicit surface with local surgery can keep resolving below it until
the mesh finally binds.

\subsubsection{And at the coarse mesh, neither method had resolved it}

\begin{table}[htbp]
\centering
\caption{The same two numbers at both resolutions, $t=4$.}
\label{tab:cuspres}
\small
\begin{adjustbox}{max width=\linewidth}
\pgfplotstabletypeset[
  col sep=space, string type,
  columns/corrida/.style={column name={}, column type={l}},
  columns/malha/.style={column name={mesh}, column type={l}},
  columns/assim/.style={column name={$A$}, column type={r}},
  columns/RRel/.style={column name={$R/R_{\mathrm{eq}}$}, column type={r}},
  columns/obs/.style={column name={}, column type={l}},
  every head row/.style={before row=\toprule, after row=\midrule},
  every last row/.style={after row=\bottomrule},
]{\dat{cuspide-resolucao.dat}}
\end{adjustbox}
\end{table}

At $h=1/25$ the two differ by a factor of \emph{fifteen} in $R/R_{\mathrm{eq}}$
--- $0.047$ against $0.709$. That is not a difference between methods; it is two
different failures. VOF at $0.709$ has a trailing edge about as round as the
bubble itself: no cusp is resolved at all. Front-tracking at $0.047$ is
\emph{spuriously} sharp, because that run carries no filter and its lowest
vertex is a zig-zag spike --- the sub-grid undulation of
Table~\ref{tab:undulation} makes the local curvature unusable. The filtered
$h=1/50$ run gives $0.086$, \emph{less} sharp than the unfiltered coarse run,
which is the direct evidence that the coarse value was undulation and not cusp.

So the paper's claim about the cusp is narrower than the one it was reaching
for, and firmer: \textbf{at the reference resolution the two methods agree on
the integrated shape of the cusp to $0.5\%$; the coarse mesh resolves it in
neither, for two unrelated reasons; and the local curvature at $t=4$ is at the
mesh floor in both and so cannot discriminate.}

\caveat{The $h=1/25$ front-tracking run is unfiltered and the $h=1/50$ one is
filtered, so the change between those two rows mixes resolution with the filter.
Separating them needs a filtered run at $h=1/25$, which the cadence sweep of
Section~\ref{sec:tsur} produces as a by-product.}

\subsection{A model limit, not a numerical one}

Oldroyd-B has no finite extensibility limit (Section~\ref{sec:formulation}), so
under sustained extension the stress grows without bound. In a preliminary run
with deliberately wrong parameters --- gravity a hundredfold too strong and
surface tension 24 times too weak --- the trace grew monotonically from 0 to 416
in 160 steps with no sign of saturation. With the correct parameters the same
configuration gives \num{3.6e-2} after seven steps, growing slowly. The
distinction matters for reading the viscoelastic results: unbounded growth is
what the model does in extension, and the question is whether it saturates on
the time scale of the rise, not whether it appears.

\subsection{Uniform mesh}

All twelve runs completed 2001 steps to $t=4$ with zero non-finite values.
That is worth stating before any number: Oldroyd-B has no finite extensibility
limit, and the question going in was whether the stress would saturate on the
time scale of the rise or run away.

\begin{table}[htbp]
\centering
\caption{The measured matrix at $t=4$, $25\times25\times75$. ``Drop'' is the
fall in rise velocity after its maximum, relative to that maximum.}
\label{tab:vemeasured}
\small
\begin{adjustbox}{max width=\linewidth}
\pgfplotstabletypeset[
  col sep=space, string type,
  columns={caso,lambda,beta,zc,wmax,tmax,west,queda},
  columns/caso/.style={column name={Run}, column type={l}},
  columns/lambda/.style={column name={$\lambda_1$}, column type={c}},
  columns/beta/.style={column name={$\beta$}, column type={c}},
  columns/zc/.style={column name={$z_c$}, column type={r}},
  columns/wmax/.style={column name={$w_{\max}$}, column type={r}},
  columns/tmax/.style={column name={at $t$}, column type={r}},
  columns/west/.style={column name={$w_{\mathrm{steady}}$}, column type={r}},
  columns/queda/.style={column name={Drop (\%)}, column type={r}},
  every head row/.style={before row=\toprule, after row=\midrule},
  every last row/.style={after row=\bottomrule},
]{\dat{ve-medido.dat}}
\end{adjustbox}
\end{table}

\subsubsection{Against the published numbers}

The reference thesis publishes curves rather than tables, but its text quotes
values, and those are what make a numerical comparison possible at all. They
come with ``approximately'' and two significant figures, which bounds how tight
the comparison can be.

\begin{table}[htbp]
\centering
\caption{Measured against the values the reference thesis states in its text.}
\label{tab:vethesis}
\small
\begin{adjustbox}{max width=\linewidth}
\pgfplotstabletypeset[
  col sep=space, string type,
  columns/caso/.style={column name={Run}, column type={l}},
  columns/zc/.style={column name={$z_c$}, column type={r}},
  columns/zcrel/.style={column name={}, column type={r}},
  columns/wmax/.style={column name={$w_{\max}$}, column type={r}},
  columns/wmaxrel/.style={column name={}, column type={r}},
  columns/west/.style={column name={$w_{\mathrm{steady}}$}, column type={r}},
  columns/westrel/.style={column name={}, column type={r}},
  every head row/.style={before row=\toprule, after row=\midrule},
  every last row/.style={after row=\bottomrule},
]{\dat{ve-contra-tese.dat}}
\end{adjustbox}
\end{table}

Every entry agrees to within 7\%, and most to within 3\%. The signs are not
random, and the next subsection returns to that once the Newtonian reference
puts nine comparisons on the table rather than six.

\subsubsection{The trends, all twelve confirmed}

The thesis makes three ordering claims: more polymer rises higher, reaches a
higher peak velocity, and settles at a higher steady velocity. Across both
methods and both relaxation times, \textbf{all twelve orderings hold}.

\begin{center}
\small
\adjustbox{max width=\linewidth}{%
\begin{tabular}{llccc}
\toprule
Method & $\lambda_1$ & $z_c$ ($\beta = 0.8, 0.5, 0.2$) & $w_{\max}$ & $w_{\mathrm{steady}}$ \\
\midrule
FT  & 0.1 & 1.8221, 1.8482, 1.8761 & 0.3642, 0.3698, 0.3754 & 0.3486, 0.3561, 0.3641 \\
FT  & 0.4 & 1.8477, 1.9156, 1.9860 & 0.3714, 0.3896, 0.4092 & 0.3557, 0.3751, 0.3972 \\
VOF & 0.1 & 1.8679, 1.8965, 1.9251 & 0.3942, 0.4002, 0.4041 & 0.3641, 0.3713, 0.3816 \\
VOF & 0.4 & 1.8942, 1.9733, 2.0540 & 0.3933, 0.4158, 0.4280 & 0.3685, 0.3948, 0.4180 \\
\bottomrule
\end{tabular}
}
\end{center}

\paragraph{One ordering appeared to fail, and the reason is methodological.}
Read before all twelve runs had reached $t=4$, the VOF series at
$\lambda_1 = 0.1$ was \emph{not} monotone in $z_c$: 1.8679, 1.8965, 1.8774. The
third run was at $t = 3.88$ while the second was at $t = 4.00$, and a centre of
mass is an integrated quantity --- comparing it across different end times
compares different things. At $t = 4.00$ the ordering holds. Every table in this
paper that reports an integrated quantity therefore also reports the end time it
was read at.

\subsubsection{The signature the thesis describes, in both methods}

The most specific claim in the reference is not an ordering but a shape: at
$\lambda_1 = 0.4$ and $\beta = 0.2$ the rise velocity falls sharply after its
maximum and then recovers. It appears in both methods --- a drop of $8.1\%$ in
front-tracking and $5.9\%$ in VOF for the cusp case, with recovery in both.

The drop is present in \emph{every} viscoelastic run here. An earlier version of
this section reported it as \emph{systematically larger in VOF}, and that claim
does not survive its own measurement --- the correction is worth stating in full,
because the failure mode is general.

\subsubsection{The asymmetry was in the measurement, not the flow}
\label{sec:dropdef}

The drop had been taken as the fall from the global maximum of $w(t)$ to the
following minimum. On a curve that is nearly flat after the initial overshoot,
that definition is governed by whichever single frame happens to sit highest ---
and in three of the seven VOF runs the ``maximum'' it selects lies at
$t \approx 3.0$--$3.9$, long after the physical overshoot at $t \approx 0.8$.

The two methods do not carry the same frame-to-frame scatter. Locating the
overshoot in the window $t \in [0.3,1.5]$ on a centred five-point mean, and
measuring the scatter of $w$ about that mean:

\begin{table}[htbp]
\centering
\caption{The same drop under two definitions, and the scatter that separates
them. ``Raw'' is global maximum then following minimum; ``robust'' locates the
overshoot physically and smooths.}
\label{tab:dropdef}
\small
\begin{adjustbox}{max width=\linewidth}
\pgfplotstabletypeset[
  col sep=space, string type,
  columns={caso,metodo,lambda,beta,qbruta,tbruta,qrobusta,trobusta,ruido},
  columns/caso/.style={column name={Run}, column type={l}},
  columns/metodo/.style={column name={}, column type={l}},
  columns/lambda/.style={column name={$\lambda_1$}, column type={r}},
  columns/beta/.style={column name={$\beta$}, column type={r}},
  columns/qbruta/.style={column name={raw (\%)}, column type={r}},
  columns/tbruta/.style={column name={$t$ of max}, column type={r}},
  columns/qrobusta/.style={column name={robust (\%)}, column type={r}},
  columns/trobusta/.style={column name={$t$ of max}, column type={r}},
  columns/ruido/.style={column name={scatter (\% of peak)}, column type={r}},
  every head row/.style={before row=\toprule, after row=\midrule},
  every last row/.style={after row=\bottomrule},
]{\dat{queda-robusta.dat}}
\end{adjustbox}
\end{table}

Under the raw definition front-tracking gives 11.5--13.8\% and VOF 11.1--23.1\%,
and VOF looks systematically worse. Under the robust definition front-tracking
gives 5.5--8.1\% (mean 6.8\%) and VOF 5.1--7.0\% (mean 6.1\%): the two
\textbf{agree}, and if anything front-tracking drops slightly more.

What differs between the methods is not the drop but the scatter --- $3.0\%$ of
the peak in front-tracking against $3.8$--$4.9\%$ in VOF, consistent with jitter
in the marching-cubes extraction from which $z_c$ is read. The raw definition
converts that scatter into an apparent physical asymmetry, because a noisier
series offers a higher stray maximum.

So the honest statement is narrower than the one it replaces: \textbf{the
velocity-drop signature the reference describes is present in both methods at
$6\pm1\%$, and these runs do not resolve a difference between them.} The
asymmetry that remains is the one in Section~\ref{sec:newtonian} --- VOF above
the published values on all nine comparisons --- which does not depend on
locating an extremum.

\subsubsection{And it does not change with $h$}

The pending test was whether the drop moves with resolution: a representation
artefact should, the physics should not. The cusp case now exists at both
resolutions, the same quantity computed the same way.

\begin{table}[htbp]
\centering
\caption{The drop at two resolutions, VOF, $\lambda_1=0.4$, $\beta=0.2$.}
\label{tab:dropres}
\small
\begin{adjustbox}{max width=\linewidth}
\pgfplotstabletypeset[
  col sep=space, string type,
  columns/malha/.style={column name={Mesh}, column type={l}},
  columns/h/.style={column name={$h$}, column type={r}},
  columns/queda/.style={column name={drop (\%)}, column type={r}},
  columns/pico/.style={column name={$w_{\max}$}, column type={r}},
  columns/tpico/.style={column name={$t_{\max}$}, column type={r}},
  columns/zcfinal/.style={column name={$z_c(4)$}, column type={r}},
  columns/psifinal/.style={column name={$\Psi(4)$}, column type={r}},
  columns/ruido/.style={column name={scatter (\%)}, column type={r}},
  every head row/.style={before row=\toprule, after row=\midrule},
  every last row/.style={after row=\bottomrule},
]{\dat{queda-resolucao.dat}}
\end{adjustbox}
\end{table}

Halving $h$ moves the drop from $5.9\%$ to $6.6\%$ --- $0.7$ percentage points,
against a frame scatter of $3$--$4\%$ of the peak. The measurement therefore
\emph{bounds} any resolution dependence below its own noise rather than
resolving one, which is consistent with the drop being physical and is not proof
of it. The rise itself converges in the usual way: $z_c(4)$ moves $0.60\%$ and
the final sphericity $0.0046$.

\subsubsection{The shapes}

Figure~\ref{fig:veshape} is the counterpart of the reference's shape figures:
the vertical section through the bubble centre, at $t=4$, for all six
$(\lambda_1,\beta)$ combinations, with the two methods overlaid. As in
Section~\ref{sec:newtonian}, contours are centred on their own centroid so that
what is compared is shape and not rise.

\begin{figure}[htbp]
\centering
\begin{tikzpicture}
\begin{axis}[width=0.325\linewidth, height=0.30\linewidth,
             axis equal image, title={$\lambda_1{=}0.1$, $\beta{=}0.8$}, grid=both, no markers,
             legend pos=south east, legend cell align=left,
             font=\scriptsize, xticklabel pos=bottom]
  \addplot+[thick] table[x=x, y=z] {\dat{secoes/ve-vof-l01-b08-c.dat}};
    \addlegendentry{VOF}
  \addplot+[thick, dashed] table[x=x, y=z] {\dat{secoes/ve-ft-l01-b08-c.dat}};
    \addlegendentry{FT}
\end{axis}
\end{tikzpicture}
\begin{tikzpicture}
\begin{axis}[width=0.325\linewidth, height=0.30\linewidth,
             axis equal image, title={$\lambda_1{=}0.1$, $\beta{=}0.5$}, grid=both, no markers,
             font=\scriptsize, xticklabel pos=bottom]
  \addplot+[thick] table[x=x, y=z] {\dat{secoes/ve-vof-l01-b05-c.dat}};
  \addplot+[thick, dashed] table[x=x, y=z] {\dat{secoes/ve-ft-l01-b05-c.dat}};
\end{axis}
\end{tikzpicture}
\begin{tikzpicture}
\begin{axis}[width=0.325\linewidth, height=0.30\linewidth,
             axis equal image, title={$\lambda_1{=}0.1$, $\beta{=}0.2$}, grid=both, no markers,
             font=\scriptsize, xticklabel pos=bottom]
  \addplot+[thick] table[x=x, y=z] {\dat{secoes/ve-vof-l01-b02-c.dat}};
  \addplot+[thick, dashed] table[x=x, y=z] {\dat{secoes/ve-ft-l01-b02-c.dat}};
\end{axis}
\end{tikzpicture}

\begin{tikzpicture}
\begin{axis}[width=0.325\linewidth, height=0.30\linewidth,
             axis equal image, title={$\lambda_1{=}0.4$, $\beta{=}0.8$}, grid=both, no markers,
             font=\scriptsize, xticklabel pos=bottom]
  \addplot+[thick] table[x=x, y=z] {\dat{secoes/ve-vof-l04-b08-c.dat}};
  \addplot+[thick, dashed] table[x=x, y=z] {\dat{secoes/ve-ft-l04-b08-c.dat}};
\end{axis}
\end{tikzpicture}
\begin{tikzpicture}
\begin{axis}[width=0.325\linewidth, height=0.30\linewidth,
             axis equal image, title={$\lambda_1{=}0.4$, $\beta{=}0.5$}, grid=both, no markers,
             font=\scriptsize, xticklabel pos=bottom]
  \addplot+[thick] table[x=x, y=z] {\dat{secoes/ve-vof-l04-b05-c.dat}};
  \addplot+[thick, dashed] table[x=x, y=z] {\dat{secoes/ve-ft-l04-b05-c.dat}};
\end{axis}
\end{tikzpicture}
\begin{tikzpicture}
\begin{axis}[width=0.325\linewidth, height=0.30\linewidth,
             axis equal image, title={$\lambda_1{=}0.4$, $\beta{=}0.2$}, grid=both, no markers,
             font=\scriptsize, xticklabel pos=bottom]
  \addplot+[thick] table[x=x, y=z] {\dat{secoes/ve-vof-l04-b02-c.dat}};
  \addplot+[thick, dashed] table[x=x, y=z] {\dat{secoes/ve-ft-l04-b02-c.dat}};
\end{axis}
\end{tikzpicture}
\caption{Viscoelastic bubble at $t=4$: vertical section through the centre, VOF
(solid) against front-tracking (dashed). Top row $\lambda_1 = 0.1$, bottom row
$\lambda_1 = 0.4$; polymer content increases to the right ($\beta$ falls).
Contours centred on their own centroid.}
\label{fig:veshape}
\end{figure}

Two things the figure shows that the tables cannot. The first is that the two
methods agree on shape across the whole matrix, including at the highest polymer
content, where the elastic stress is largest. The second is the trailing edge:
at $\lambda_1 = 0.4$, $\beta = 0.2$ --- the bottom-right panel --- the rear of
the bubble is drawn out relative to every other case, which is the cusp the
reference describes.

\paragraph{What this figure does not settle.} Whether the drawn-out rear is
\emph{resolved} or merely \emph{suggested} at this mesh. A cusp is curvature
concentrated on a scale that approaches the cell size, and this is the coarse
mesh. \missing{the reference resolution for this case is running; the comparison
of the two will say whether the feature sharpens with $h$ --- which it should, if
it is the cusp --- or stays at the same width, which would mean the mesh is
setting its scale}

\subsubsection{The Newtonian reference, in the same geometry}

One claim of the reference could not be tested until now: that the viscoelastic
bubble rises higher than a Newtonian one. The Newtonian runs of
Section~\ref{sec:newtonian} cannot serve as that reference --- they use a
different box and end time, and the centre of mass is an \emph{integrated}
quantity, so it is not comparable across either. A Newtonian pair was therefore
run in the thesis geometry: box $1\times1\times3$, to $t=4$.

\begin{table}[htbp]
\centering
\caption{The trend claim, with a reference in the same geometry. Centre of mass
at $t=4$.}
\label{tab:vetendency}
\small
\pgfplotstabletypeset[
  col sep=space, string type,
  columns/caso/.style={column name={Case ($\beta=0.2$)}, column type={l}},
  columns/ve/.style={column name={viscoelastic}, column type={r}},
  columns/newt/.style={column name={Newtonian}, column type={r}},
  columns/veredito/.style={column name={}, column type={l}},
  every head row/.style={before row=\toprule, after row=\midrule},
  every last row/.style={after row=\bottomrule},
]{\dat{ve-tendencia.dat}}
\end{table}

All four hold. And the Newtonian reference itself lands close to the published
value: front-tracking gives $z_c = 1.8055$ against $\approx 1.8$, and a peak
velocity of $0.3603$ against $\approx 0.36$ --- $0.3\%$ and $0.1\%$.

\paragraph{The asymmetry in the signs is now worth stating as a result.}
With the Newtonian reference included there are nine comparisons against
published values per method. \textbf{VOF is above the reference on all nine};
front-tracking is at or below on seven, and its two exceptions are $+0.1\%$ and
$+0.3\%$. The mean absolute deviation is $1.2\%$ for front-tracking against
$3.6\%$ for VOF.

That is no longer the scatter of two or three readings, and it is the opposite
of what the provenance would suggest: the reference was computed with VOF. Two
readings are available and these runs do not separate them --- our
volume-of-fluid differs from theirs in some detail of the reconstruction or the
initial fraction, or the marching-cubes extraction from which our shape
quantities are read is itself biased upward. \missing{the second is testable:
the field-based centroid probe agrees with the surface-based one to
\num{1e-4} in the Newtonian case of Section~\ref{sec:newtonian}, which argues
against extraction bias, but that check has not been repeated here}

\subsection{Adaptive mesh}
\missing{blocked on Section~\ref{sec:amr}}

\section{Cost and performance}
\label{sec:cost}

\status{Uniform-mesh rates are \textbf{measured} at both resolutions. The
refinement gain is bounded by two measured numbers --- cell count and remeshing
overhead --- but the production comparison has not been run.}

\subsection{Uniform mesh}

\begin{table}[htbp]
\centering
\caption{Measured cost in microseconds per cell-step, single core, uniform mesh.
Same mesh and same conditions for both methods.}
\label{tab:cost}
\small
\pgfplotstabletypeset[
  col sep=space, string type,
  columns/caso/.style={column name={Rheology}, column type={l}},
  columns/vof/.style={column name={VOF}, column type={r}},
  columns/ft/.style={column name={FT}, column type={r}},
  every head row/.style={before row=\toprule, after row=\midrule},
  every last row/.style={after row=\bottomrule},
]{\dat{custo.dat}}
\end{table}

The two methods cost the same to within five percent, and which is cheaper
depends on the resolution: front-tracking is 2.5\% more expensive at
$25\times25\times50$ and 4.9\% cheaper at $50\times50\times100$ (19.5 against
20.5 hours of wall clock for the full 1501 steps).

The weak version of that statement is the defensible one, and it is still worth
making, because the strong opposite is often assumed: an explicit surface sounds
much more expensive than a scalar field, and it is not. VOF pays for the PLIC
reconstruction and the split advection in every cell near the interface;
front-tracking pays for advecting a few thousand vertices and rebuilding one
spatial index. Those turn out to be comparable.

\paragraph{How not to measure this.}
An earlier version of this table reported front-tracking as 12\% cheaper. That
figure came from eleven-step runs, where setup is a large fraction of the total
--- and, it turned out, a \emph{different} fraction for the two methods. Only
the complete runs give a usable ratio.

\paragraph{The cost per cell-step rises with problem size.}
By 14\% for VOF and 6\% for front-tracking between the two resolutions. Part of
that is cache pressure and part is that the two reference runs shared the
machine on separate cores while the coarse runs did not; these measurements do
not separate the two causes. Every rate in this paper states the problem size
and whether the machine was shared, because a rate without those is not
transferable --- an early estimate in this work extrapolated a small-problem
rate to the large problem and underpredicted the wall clock by a factor of about
1.6.

\subsection{Output cost}

Worth a paragraph because it dominated a first attempt at the reference
resolution. The field output is an unstructured grid with eight unshared corner
points per cell: $250\,000$ cells become two million points, about 300~MB per
frame, roughly 20 minutes to write. At 120 frames that is some 40 hours of
writing against 16 hours of computation. The extracted interface is a factor of
226 smaller and already carries the interpolated fields at its vertices, so the
field output is written at a coarse cadence and the interface at every frame.

\subsection{Adaptive against uniform}

The remeshing machinery works (Section~\ref{sec:amr}); what is missing is a
production run on an adapted mesh. Two measured numbers already bound what to
expect.

\paragraph{Cell count.} For the reference geometry with a minimum of five fine
cells per side, the criterion produces, against the equivalent uniform fine
mesh:

\begin{center}
\small
\begin{tabular}{llrrr}
\toprule
Base & Levels & Criterion & Uniform fine & Factor \\
\midrule
$50\times50\times100$ & 1 & 379\,136 & 2\,000\,000 & 5.3 \\
$50\times50\times100$ & 2 & 1\,128\,584 & 16\,000\,000 & 14.2 \\
$50\times50\times100$ & 3 & 5\,458\,224 & 128\,000\,000 & 23.5 \\
\bottomrule
\end{tabular}
\end{center}

Base 50 with two levels is an effective $200\times200\times400$ at 1.13 million
cells. At the measured rates that is of order five days against roughly seventy
for the equivalent uniform mesh --- which is what makes the fine resolution
reachable at all.

\paragraph{Remeshing overhead.} The criterion costs 0.008~s; the reconstruction
costs between 1.9 and 4.7~s, varying with how much the mesh changed. At one
remesh every 25 steps that is under 1\% of a step. The bottleneck is the
reconstruction, not the criterion, by a factor of 230.

\paragraph{The band over-delivers, and it costs cells.} The measured band is
8.0, 12.1 and 16.1 fine cells for one, two and three levels, against a floor of
five. The margin in the rule is fixed in \emph{base}-cell units, so it is worth
$2^{\text{levels}}$ fine cells. The rule was calibrated by measurement in two
dimensions --- half the margin undershot the floor --- so it is carried over
unchanged here rather than diverging between the two codes, but a margin
expressed in fine cells would recover a factor of two or three in cell count at
three levels.

\missing{the production comparison: same accuracy for how much less}

\section{Conclusions}
\label{sec:conclusions}

\status{Depends on the sections above. What can be written now is listed; the
rest is marked.}

\missing{conclusions --- to be written once the matrix is closed}

\subsection*{What is established}

\begin{itemize}
\item Both methods reproduce the ellipsoidal regime of the Newtonian reference
case and the shape of its published rise-velocity curve, and they \emph{converge
toward each other} under refinement --- from 1.94\% to 0.60\% in centre of mass
and from 4.3\% to 1.1\% in terminal rise velocity --- approaching from opposite
sides.
\item At the reference resolution \textbf{VOF} conserves mass about five times
better ($-0.170\%$ against $-0.851\%$), and front-tracking's drift \emph{grows}
with refinement. The coarse-mesh figures reverse this ranking and are partly
cancellation; they should not be quoted.
\item The two methods cost the same to within 5\% per cell-step, with the sign
of the difference depending on resolution.
\item The verification apparatus separates exact from convergent oracles, and
the exact ones --- surface closure, genus, the plane/slab/tetrahedron fraction
--- pass at machine precision.
\item The viscoelastic matrix reproduces the reference results. All twelve
ordering claims of the reference thesis hold, the quoted values agree to within
7\% (most within 3\%), the trend against a Newtonian reference in the same
geometry holds in all four cases, and the specific velocity signature the
reference describes --- a sharp drop after the maximum, then recovery ---
appears in both methods.
\item \textbf{VOF sits above the published values on all nine comparisons and
front-tracking at or below on seven of nine}, with mean absolute deviations of
$3.6\%$ and $1.2\%$. The asymmetry is the opposite of what the provenance
suggests, since the reference was computed with VOF.
\item Oldroyd-B did \emph{not} run away on the time scale of the rise: all
twelve runs completed 2001 steps with zero non-finite values.
\item \textbf{The cusp pair is closed at the reference resolution, and the two
methods agree on the integrated shape of the cusp to $0.5\%$} --- asymmetry
$1.2846$ in front-tracking against $1.2907$ in VOF, tracking each other to
$0.2$--$2.4\%$ from $t=1.6$ onward. This is the discriminating case, with the
two interface representations at their most different.
\item But their agreement in the \emph{local} curvature at $t=4$ is not
convergence: both give $R \approx 1.07\,h$, one cell. Two methods pinned at the
same mesh floor agree for a reason unrelated to the physics. What the column
does show is the \emph{approach}: VOF reaches its floor by $t=2.4$ and sits
there, front-tracking descends gradually and arrives only at $t=4$ --- the
signature of a per-cell planar reconstruction saturating at the cell while an
explicit surface keeps resolving below it.
\item At $h=1/25$ neither method had resolved the cusp, for two unrelated
reasons: VOF reports a trailing edge as round as the bubble ($R/R_{\mathrm{eq}}
= 0.71$), and the unfiltered front-tracking reports a spuriously sharp one
($0.047$) because its lowest vertex is a zig-zag spike. The filtered $h=1/50$
run gives $0.086$ --- \emph{less} sharp than the unfiltered coarse run, which is
the direct evidence that the coarse value was undulation.
\item The drop is $5.9\%$ at $h=1/25$ and $6.6\%$ at $h=1/50$ --- a bound below
the measurement's own noise rather than a resolved dependence.
\end{itemize}

\subsection*{What was withdrawn}

One result in this paper was retracted by its own measurement, and it is recorded
here rather than quietly dropped. An earlier version reported the velocity drop
as \emph{systematically larger in VOF} ($8.6$--$13.5\%$) than in front-tracking
($4.1$--$6.6\%$), and offered it as a difference between the two interface
representations. That number came from taking the global maximum of $w(t)$ on a
curve that is nearly flat after the overshoot --- so it was set by whichever
single frame sat highest, and in three of seven VOF runs it selected a frame at
$t \approx 3$--$4$, long after the physical overshoot. Locating the overshoot
physically and smoothing, front-tracking gives $5.5$--$8.1\%$ and VOF
$5.1$--$7.0\%$: the methods agree. What actually differs is the frame-to-frame
scatter, $3.0\%$ of the peak against $3.8$--$4.9\%$, and the raw definition
converted that scatter into an apparent physical asymmetry
(Section~\ref{sec:dropdef}).

The scatter is itself evidence on an open question below: a noisier extraction is
what the marching-cubes-bias reading predicts.

\subsection*{What is not}

\begin{itemize}
\item Adaptive refinement in production. The dynamic remeshing loop works and
is verified, but no case has been run on an adapted mesh to completion.
\item \emph{Resolved, and moved here from the open list.} Front-tracking's
mass drift appeared to grow with refinement because both production runs held
the surface subdivision fixed while the mesh was refined, so $ds/h$ doubled
between them. With $ds/h$ held constant the drift falls at an observed order of
$1.87$, and surgery reduces it by $42\%$ rather than causing it
(Section~\ref{sec:deriva-nsub}). What remains open is narrower: whether the
same holds at the production resolutions, which requires rerunning the coarse
case with the surface resolution matched to the grid.
\item Whether the $h=1/50$ agreement on the cusp survives separating resolution
from the filter. The coarse front-tracking run is unfiltered and the fine one is
filtered, so those two rows differ in two things at once. The cadence sweep of
Section~\ref{sec:tsur} produces a filtered coarse run as a by-product and
settles it.
\item Which of the two readings explains the sign asymmetry: a difference
between our volume-of-fluid and the reference's, or an upward bias in the
marching-cubes extraction from which the shape quantities are read.
\item Case~2, where topological change is expected.
\item Any statement about more than one process.
\end{itemize}


% \nocite{*} faz a bibliografia do esqueleto listar TODAS as entradas
% disponiveis em referencias.bib, inclusive as ainda nao citadas no texto.
% REMOVER NA SUBMISSAO -- sem isso, uma secao sem \cite sai com bibliografia
% vazia e o PDF compila sem reclamar.
\nocite{*}

\bibliographystyle{elsarticle-num}
\bibliography{referencias}

\end{document}
